% >>> typage-template acl@2026.10.1 variant=review
\documentclass[11pt]{article}
\usepackage[review]{acl}
\usepackage{times}
\usepackage{latexsym}
\usepackage[T1]{fontenc}
\usepackage{inconsolata}
% <<< typage-template
\usepackage{microtype}
\usepackage{graphicx}
\usepackage{booktabs}
\usepackage{amsmath}

\title{Paper Title}

\author{First Author \\
  Department, University \\
  City, Country \\
  \texttt{first.author@example.com} \\\And
  Second Author \\
  Institute \\
  City, Country \\
  \texttt{second.author@example.com} \\}

\begin{document}
\maketitle

\begin{abstract}
State the problem, the approach and the main result in a single
paragraph.
\end{abstract}

\section{Introduction}
Introduce the problem and list the contributions. Cite prior work in
running text, as in \citet{knuth1984texbook}, or in
parentheses~\citep{lamport1994latex}.

\section{Method}
Describe the method. Number the equations you refer to:
\begin{equation}
  \label{eq:objective}
  \min_{\theta} \; \frac{1}{n} \sum_{i=1}^{n} \ell\bigl(f_\theta(x_i), y_i\bigr) .
\end{equation}

\section{Experiments}
Table~\ref{tab:results} reports the results of
optimizing~\eqref{eq:objective}. Describe the data, the baselines and the
compute used, and report the variation across runs.

\begin{table}[!htb]
  \caption{Results on the example benchmark (mean $\pm$ standard deviation over five runs).}
  \label{tab:results}
  \centering
  \begin{tabular}{lcc}
    \toprule
    Method   & Accuracy (\%)  & Time (s) \\
    \midrule
    Baseline & 81.2 $\pm$ 0.4 & 12.4     \\
    Ours     & 84.7 $\pm$ 0.3 & 10.9     \\
    \bottomrule
  \end{tabular}
\end{table}

\section{Conclusion}
Conclude and outline future work.

\section*{Limitations}
Describe the limitations of the work. ACL Rolling Review requires this
section; it does not count towards the page limit.

\bibliography{refs}

\appendix
\section{Example Appendix}
Optional appendices go after the references.

\end{document}
